\begin{papercard}{braun2006}
\cf{Stages} \stagetag{C} (also \stagetag{I} \stagetag{R}). The foundational paper on thematic analysis (TA): a definition, the decisions that must be made explicit, a six-phase procedure, pitfalls, and a 15-point checklist. \emph{Note:} the appended file is the \emph{UWE accepted manuscript}, not the journal typesetting (printed pp.\ 77--101), so page numbers below refer to the manuscript. Read it together with Braun \& Clarke (2021), which corrects widespread misreadings.

\cf{Positioning} TA is ``a method for identifying, analysing, and reporting patterns (themes) within data'' that minimally organises and describes the data set in rich detail and often interprets it \pc{braun2006}{6}{Thematic analysis is a method for identifying, analysing, and reporting patterns (themes) within}. It is a ``poorly branded'' method, so the authors argue it should be a method in its own right, independent of any one theory, and distinct from IPA, grounded theory, thematic discourse analysis and thematic decomposition \pc{braun2006}{8}{Thematic analysis differs from other analytic methods that seek to describe patterns across qualitative data}. Aim: to demarcate it clearly while keeping its flexibility, making researchers state assumptions and the often-omitted ``how'' \pc{braun2006}{5}{Qualitative psychologists need to be clear about what they are doing and why, and include the often-omitted}. Terms: \emph{data corpus} (all data), \emph{data set} (the part used), \emph{data item}, \emph{data extract} \pc{braun2006}{5}{Data corpus refers to all data collected for a particular research project, while data set refers to all the data}.

\cf{Decisions to make before analysis} (1) \textbf{What counts as a theme}: it captures something important in relation to the research question and shows some patterned response, but there is no fixed prevalence threshold (50\,\% vs.\ 47\,\% is irrelevant); prevalence can be counted per item, per speaker or per occurrence, and must be applied consistently \pc{braun2006}{10}{It is not the case that if it was present in 50\% of one’s data items, it would be a theme}. (2) \textbf{Rich description of the whole data set vs.\ detailed account of one aspect}. (3) \textbf{Inductive vs.\ theoretical (deductive)}: inductive codes without fitting a pre-existing frame, theoretical is analyst-driven and less rich overall; no analysis is free of theory \pc{braun2006}{12}{Inductive analysis is therefore a process of coding the data without trying to fit it into a pre-existing coding frame}. (4) \textbf{Semantic vs.\ latent} (surface vs.\ underlying assumptions; the ``jelly'' analogy) \pc{braun2006}{13}{With a semantic approach, the themes are identified within the explicit or surface meanings of the data}. (5) \textbf{Epistemology}: realist/essentialist vs.\ constructionist, with a contextualist middle \pc{braun2006}{14}{The research epistemology guides what you can say about your data, and informs how you theorise meaning}. Interview questions are not the themes \pc{braun2006}{15}{Some of the worst examples of ‘thematic’ analysis we have read have simply used the questions put to participants as the ‘themes’}.

\cf{Phase 1: familiarise} Repeated, active reading of the entire data set before coding; take notes; transcribe if needed (orthographic/verbatim minimum, checked against recordings) \pc{braun2006}{16}{It is ideal to read through the entire data set at least once before you begin your coding}. Writing starts in phase 1. Timing of the literature depends on inductive vs.\ theoretical orientation \pc{braun2006}{16}{a more inductive approach would be enhanced by not engaging with literature in the early stages of analysis}.

\cf{Phase 2: initial codes} Codes mark an interesting feature (semantic or latent), the most basic segment of data that can be assessed meaningfully; they differ from broader themes. Work through the whole data set with equal attention. Advice: code for as many potential themes as possible, code \emph{inclusively} (keep surrounding context), an extract can be uncoded, coded once or many times, and keep deviant accounts \pc{braun2006}{19}{Key advice for this phase is: a) code for as many potential themes/patterns as possible}.

\cf{Phase 3: searching for themes} Sort codes into candidate themes and sub-themes with tables, mind-maps or paper piles. A temporary ``miscellaneous'' theme is acceptable, and nothing is abandoned yet \pc{braun2006}{20}{it is perfectly acceptable to create a ‘theme’ called miscellaneous to house the codes}. \textbf{Phase 4: reviewing}: level 1 against coded extracts, level 2 against the whole data set; apply Patton's internal homogeneity and external heterogeneity; re-code as needed but \emph{stop} when refinements add nothing (the cake-decorating analogy) \pc{braun2006}{21}{when your refinements are not adding anything substantial, stop}. \textbf{Phase 5: defining and naming}: capture the essence of each theme and its scope in a couple of sentences, avoid overloaded themes, sub-themes as themes-within-a-theme, concise names \pc{braun2006}{22}{One test for this is to see whether you can describe the scope and content of each theme in a couple of sentences}. \textbf{Phase 6: report}: a concise, coherent, non-repetitive narrative with vivid extracts embedded in an argument that answers the research question \pc{braun2006}{23}{Extracts need to be embedded within an analytic narrative}. The process is recursive, not linear, and should not be rushed \pc{braun2006}{16}{Moreover, analysis is not a linear process where you simply move from one phase to the next}.

\cf{Interpretation and pitfalls} Analytic claims must go beyond the surface: what does the theme mean, what assumptions underlie it, what are its implications, why this way of talking \pc{braun2006}{24}{what does this theme mean? ’ ‘What are the assumptions underpinning it?}. Pitfalls: (1) no real analysis (strung-together or paraphrased extracts); (2) using the interview questions as themes; (3) weak analysis with overlapping or incoherent themes or only one or two extracts (``anecdotalism''); (4) mismatch between data and claims; (5) mismatch between theory and claims; and failing to spell out assumptions \pc{braun2006}{25}{The first of these is a failure to actually analyse the data at all}.

\cf{Quality (Table 2, 15 points)} Transcription (1); coding (2--6: equal attention, thorough rather than anecdotal, all extracts collated, themes checked against each other and the data, themes coherent and distinctive); analysis (7--10: interpreted not paraphrased, claims match extracts, convincing story, balance of narrative and extracts); overall (11: enough time); report (12--15: assumptions explicit, described method matches reported analysis, language consistent with epistemology, researcher positioned as active so that themes do not just ``emerge'') \pc{braun2006}{36}{A 15-Point Checklist of Criteria for Good Thematic Analysis}. ``Rigour lies in devising a systematic method whose assumptions are congruent with the way one conceptualises the subject matter'' \pc{braun2006}{27}{rigour lies in devising a systematic method whose assumptions are congruent}.

\cf{Advantages and limits} Flexible, accessible, summarises large data sets, highlights similarities and differences, suited to policy work. Limits: little interpretive power without a theoretical anchor, loses within-person continuity, no claims about language use \pc{braun2006}{27}{a thematic analysis has limited interpretative power beyond mere description if it is not used within an existing theoretical framework}.

\cf{Use for this thesis} Use it for vocabulary (data corpus/set/item/extract; code vs.\ theme; semantic vs.\ latent; inductive vs.\ theoretical) and for the reporting items (state assumptions, describe what was actually done, avoid ``themes emerged''). It does \emph{not} support a fixed codebook, independent coders or reliability statistics (Braun \& Clarke 2021 explicitly say so), so do not cite it for those. The walkthrough is a semantic, theory-informed (deductive) analysis. State that, and treat the six phases as a generic outline of the qualitative synthesis step, not as the coding-reliability design.
\end{papercard}
